\makeatletter
\def\input@path{{./}{../}{../../}{../../../}{../../../../}{preamble/}{../preamble/}{../../preamble/}{../../../preamble/}}
\makeatother
\input{preamble_exams.tex}
\title{Grade 11 -- Term 1 \\ C8: Student's t in Context}
\author{}\date{}
\pagestyle{empty}
\linespread{1.05}
\setlength{\parskip}{3pt}
\setlength{\columnsep}{7mm}
\AtBeginDocument{\setlength{\abovedisplayskip}{4pt}\setlength{\belowdisplayskip}{4pt}}
\begin{document}
	\newgeometry{left=0.4in,right=0.4in,top=0.35in,bottom=0.35in}
	\begin{center}{\large\bfseries G11 T1 C8 -- Student's t: Classroom Activity}\end{center}
	\textbf{Models and simplifying approximation (all pages).} Each question supplies a continuous model for a signed deviation: $X=m+aT$, where $T\sim t_\nu$, $m$ is the location parameter (and mean for $\nu>1$), $a>0$ is the scale, and $\boxed{\nu=n-1}$, where $n$ is the number of observations in the sample. For simplicity, the quantity described as an ``approximate standard deviation'' is used directly as the scale $a$. This is a teaching approximation: the model's actual standard deviation is $\operatorname{SD}(X)=a\sqrt{\frac{\nu}{\nu-2}}$ for $\nu>2$, which is larger than $a$. The approximation is coarser for small $\nu$ and improves as $\nu$ increases. Use the supplied model without adjusting $a$. Negative deviations mean below the reference; positive deviations mean above it.
	
	\textbf{Using the supplied table.} In row $\nu$ and column $\alpha$, the positive critical value $t^*$ satisfies $P(T>t^*)=\alpha$. By symmetry, $P(T<-t^*)=\alpha$. Consequently, $P(T\le t^*)=1-\alpha$ and $P(T\le -t^*)=\alpha$. Standardize using $t=\frac{x-m}{a}$; back-transform using $x=m+at$. These calculations use the supplied model and its scale $a$. Use only table entries (no CDF calculator). Table values are rounded to three decimals, so cutoffs computed from them are approximate. Numerical bounds in direct questions retain enough digits to recover the intended table entries. Report probabilities exactly as implied by the table and contextual cutoffs to two decimals. Endpoints have probability zero.
	
	\begin{center}\textbf{Worked Examples}\end{center}
	\begin{enumerate}[leftmargin=*,itemsep=7pt,parsep=0pt,topsep=2pt]
		\item \textbf{Direct left tail -- arrival error.} The delivery-time deviation from schedule of 8 deliveries, in minutes, has mean $1$ and approximate standard deviation $4$. Model $X=1+4T$, where $T\sim t_{7}$. Find the probability that the deviation is below $-6.58$, i.e., $P(X<-6.58)$.
		\par Standardize: $t=\frac{-6.58-1}{4}=-1.895$. Row $7$, $\alpha=0.05$: $P(T<-1.895)=\boxed{0.05}$. About $5\%$ of deliveries arrive more than $6.58$ minutes early.
		
		\item \textbf{Direct right tail -- cost error.} The cost deviation from budget of 12 batches, in dollars, has mean $-10$ and approximate standard deviation $20$. Model $X=-10+20T$, where $T\sim t_{11}$. Find the probability that the deviation exceeds $17.26$, i.e., $P(X>17.26)$.
		\par Standardize: $t=\frac{17.26+10}{20}=1.363$. Row $11$, $\alpha=0.10$: $P(T>1.363)=\boxed{0.10}$. About $10\%$ of batches exceed budget by more than $\$17.26$.
		
		\item \textbf{Direct interval -- energy error.} The energy-use deviation from forecast over 15 days, in kWh, has mean $5$ and approximate standard deviation $10$. Model $X=5+10T$, where $T\sim t_{14}$. Find the probability that the deviation lies between $-3.68$ and $22.61$, i.e., $P(-3.68<X<22.61)$.
		\par Standardize: $t_L=\frac{-3.68-5}{10}=-0.868$ and $t_U=\frac{22.61-5}{10}=1.761$. Row $14$: $P(T<-0.868)=0.20$, $P(T>1.761)=0.05$. Thus $P(-0.868<T<1.761)=1-0.20-0.05=\boxed{0.75}$.
		
		\item \textbf{Direct combined tails -- revenue alerts.} The revenue deviation from plan over 6 weeks, in thousands of dollars, has mean $0$ and approximate standard deviation $2$. Model $X=2T$, where $T\sim t_{5}$. An alert occurs when the deviation is below $-5.142$ or above $2.952$. Find the alert probability, i.e., $P(X<-5.142\text{ or }X>2.952)$.
		\par Standardize: $t_L=\frac{-5.142-0}{2}=-2.571$ and $t_U=\frac{2.952-0}{2}=1.476$. Row $5$ gives lower tail $0.025$ and upper tail $0.10$. Disjoint tails add: $\boxed{0.125}$, or $12.5\%$.
		
		\item \textbf{Inverse left tail -- temperature error.} The reading deviation from reference of 10 sensors, in degrees Celsius, has mean $2$ and approximate standard deviation $4$. Model $X=2+4T$, where $T\sim t_{9}$. Find the cutoff $x$ below which $10\%$ of deviations fall, i.e., $P(X<x)=0.10$.
		\par Row $9$, $\alpha=0.10$: $t=-1.383$. Hence $x=2+4(-1.383)\approx\boxed{-3.53\ ^\circ\mathrm{C}}$. About $10\%$ of errors fall below this cutoff.
		
		\item \textbf{Inverse right tail -- processing error.} The processing-time deviation from target of 18 batches, in minutes, has mean $-1$ and approximate standard deviation $3$. Model $X=-1+3T$, where $T\sim t_{17}$. Find the cutoff $x$ exceeded by $5\%$ of deviations, i.e., $P(X>x)=0.05$.
		\par Row $17$, $\alpha=0.05$: $t=1.740$. Hence $x=-1+3(1.740)\approx\boxed{4.22\text{ minutes}}$.
		
		\item \textbf{Inverse interval -- demand error.} The demand deviation from forecast over 14 weeks, in tonnes, has mean $5$ and approximate standard deviation $10$. Model $X=5+10T$, where $T\sim t_{13}$. Find limits $L<U$ enclosing the central $90\%$ of deviations, i.e., $P(L<X<U)=0.90$, with equal excluded tails.
		\par Each tail is $0.05$, so row $13$ gives $t^*=1.771$. Thus $(L,U)=(5-10(1.771),5+10(1.771))\approx\boxed{(-12.71,22.71)\text{ tonnes}}$.
		
		\item \textbf{Inverse combined tails -- filling alerts.} The fill-mass deviation from target of 5 containers, in grams, has mean $0$ and approximate standard deviation $2$. Model $X=2T$, where $T\sim t_{4}$. Find limits $L<U$ leaving $5\%$ of deviations outside, equally split between the tails, i.e., $P(X<L)=P(X>U)=0.025$.
		\par Each tail is $0.025$. Row $4$: $t^*=2.776$. Therefore $(L,U)=(-2(2.776),2(2.776))\approx\boxed{(-5.55,5.55)\text{ g}}$; $2.5\%$ lie in each tail.
	\end{enumerate}
	
	
	\par\textbf{PROBLEMS}\par
	\par\textbf{Instructions.} Use only the supplied Student's $t$ critical-value table. Every specified standardized boundary is a rounded table entry. For direct problems, standardize and identify the appropriate tail probability; for inverse problems, find the table critical value and back-transform. Show your calculations and interpret results in context. No CDF or inverse-CDF calculator is needed. Show working on separate paper.
	\begin{enumerate}[leftmargin=*,itemsep=12pt,parsep=0pt,topsep=6pt]
		\item \textbf{Bus arrivals.} The arrival-time deviation from schedule of 7 buses, in minutes, has mean $2$ and approximate standard deviation $6$. Model $X=2+6T$, where $T\sim t_{6}$. Find the probability that the deviation is below $-9.658$, i.e., $P(X<-9.658)$, and interpret the threshold.
		\item \textbf{Shipping costs.} The cost deviation of 16 shipments, in thousand USD, has mean $-15$ and approximate standard deviation $40$. Model $X=-15+40T$, where $T\sim t_{15}$. Find the probability that the deviation exceeds $-49.64$, i.e., $P(X>-49.64)$, and explain why it exceeds $0.5$.
		\item \textbf{Water demand.} The water-use deviation from forecast over 20 days, in cubic metres, has mean $10$ and approximate standard deviation $20$. Model $X=10+20T$, where $T\sim t_{19}$. Find the probability that the deviation lies between $27.22$ and $44.58$, i.e., $P(27.22<X<44.58)$.
		\item \textbf{Harvest forecasts.} The harvested-mass deviation from forecast of 9 harvests, in tonnes, has mean $-2$ and approximate standard deviation $4$. Model $X=-2+4T$, where $T\sim t_{8}$. Find the probability that the deviation lies between $-11.224$ and $-5.556$, i.e., $P(-11.224<X<-5.556)$.
		\item \textbf{Weekly revenue alerts.} The revenue deviation from plan over 5 weeks, in thousands of dollars, has mean $1$ and approximate standard deviation $3$. Model $X=1+3T$, where $T\sim t_{4}$. An alert occurs when the deviation is below $-5.396$ or above $5.599$. Find the alert probability, i.e., $P(X<-5.396\text{ or }X>5.599)$.
		\item \textbf{Cold-storage sensors.} The reading deviation from reference of 11 sensors, in degrees Celsius, has mean $-2$ and approximate standard deviation $6$. Model $X=-2+6T$, where $T\sim t_{10}$. Find the cutoff $x$ below which $10\%$ of deviations fall, i.e., $P(X<x)=0.10$.
		\item \textbf{Rail journey delays.} The journey-time deviation from schedule of 24 rail journeys, in minutes, has mean $3$ and approximate standard deviation $8$. Model $X=3+8T$, where $T\sim t_{23}$. Find the cutoff $x$ exceeded by only $2.5\%$ of journeys, i.e., $P(X>x)=0.025$.
		\item \textbf{Electricity forecasting.} The electricity-use deviation from forecast over 19 weeks, in kWh, has mean $10$ and approximate standard deviation $40$. Model $X=10+40T$, where $T\sim t_{18}$. Find limits $L<U$ enclosing the central $90\%$ of deviations, i.e., $P(L<X<U)=0.90$, with equal excluded tails.
		\item \textbf{Packaging inspection.} The mass deviation from target of 6 packages, in grams, has mean $-0.5$ and approximate standard deviation $1.5$. Model $X=-0.5+1.5T$, where $T\sim t_{5}$. Find limits $L<U$ leaving $5\%$ of deviations outside, equally split between the tails, i.e., $P(X<L)=P(X>U)=0.025$.
		\item \textbf{Processing-time alerts.} The processing-time deviation from target of 13 batches, in minutes, has mean $1$ and approximate standard deviation $5$. Model $X=1+5T$, where $T\sim t_{12}$. Find limits $L<U$ leaving $15\%$ of deviations outside, with $5\%$ below $L$ and $10\%$ above $U$, i.e., $P(X<L)=0.05$ and $P(X>U)=0.10$.
	\end{enumerate}
	
	\newpage
	\begin{center}{\large\bfseries Detailed Solutions}\end{center}
	\small\linespread{1}\selectfont
	\setlength{\abovedisplayskip}{2pt}\setlength{\belowdisplayskip}{2pt}
	\setlength{\jot}{1pt}
	All critical values are read from the supplied table and rounded to three decimals. Probabilities are table-implied; transformed cutoffs are approximate.
	\begin{multicols}{2}
		\begin{enumerate}[leftmargin=*,itemsep=5pt,parsep=1pt,topsep=2pt]
			\item \textbf{Bus arrivals ($\nu=6$).}
			\par $t=\frac{-9.658-2}{6}=-1.943$. Row $6$, $\alpha=0.05$:
			\[P(T<-1.943)=\boxed{0.05}.\]
			The cutoff is approximately $-9.66$ minutes; about $5\%$ arrive more than $9.66$ minutes early.
			
			\item \textbf{Shipping costs ($\nu=15$).}
			\par $t=\frac{-49.64+15}{40}=-0.866$. Row $15$, $\alpha=0.20$:
			\[P(T>-0.866)=1-0.20=\boxed{0.80}.\]
			The threshold is $-49.64$ thousand USD relative to budget. It is below the location $-15$ thousand USD, so more than half of outcomes exceed it.
			
			\item \textbf{Water demand ($\nu=19$).}
			\par $t_L=\frac{27.22-10}{20}=0.861$, $t_U=\frac{44.58-10}{20}=1.729$. Row $19$: $P(T>0.861)=0.20$ and $P(T>1.729)=0.05$.
			\[P(0.861<T<1.729)=0.20-0.05=\boxed{0.15}.\]
			The bounds are $27.22$ and $44.58$ cubic metres above forecast; about $15\%$ of days fall between them.
			
			\item \textbf{Harvest forecasts ($\nu=8$).}
			\par $t_L=\frac{-11.224+2}{4}=-2.306$, $t_U=\frac{-5.556+2}{4}=-0.889$. Row $8$ gives $P(T<-2.306)=0.025$ and $P(T<-0.889)=0.20$.
			\[P(-2.306<T<-0.889)=0.20-0.025=\boxed{0.175}.\]
			The interval is approximately $(-11.22,-5.56)$ tonnes; about $17.5\%$ of harvests fall short by this amount.
			
			\item \textbf{Revenue alerts ($\nu=4$).}
			\par $t_L=\frac{-5.396-1}{3}=-2.132$, $t_U=\frac{5.599-1}{3}=1.533$. Row $4$: lower tail $0.05$, upper tail $0.10$.
			\[P(\text{alert})=0.05+0.10=\boxed{0.15}.\]
			The alerts occur below $-5.40$ or above $5.60$ thousand dollars relative to plan. About $15\%$ of weeks trigger an alert.
			
			\columnbreak
			\item \textbf{Cold-storage sensors ($\nu=10$).}
			\par Lower-tail probability $0.10$ corresponds to $t=-1.372$ (row $10$).
			\[x=-2+6(-1.372)\approx\boxed{-10.23\ ^\circ\mathrm{C}}.\]
			About $10\%$ of sensor errors fall below $-10.23\ ^\circ\mathrm{C}$ relative to reference.
			
			\item \textbf{Rail journey delays ($\nu=23$).}
			\par Upper-tail probability $0.025$ corresponds to $t=2.069$ (row $23$).
			\[x=3+8(2.069)\approx\boxed{19.55\text{ minutes}}.\]
			About $2.5\%$ of journeys are more than $19.55$ minutes late.
			
			\item \textbf{Electricity forecasting ($\nu=18$).}
			\par Central $90\%$ leaves $0.05$ in each tail. Row $18$, $\alpha=0.05$: $t^*=1.734$.
			\[L=10-40(1.734)\approx\boxed{-59.36\text{ kWh}},\]
			\[U=10+40(1.734)\approx\boxed{79.36\text{ kWh}}.\]
			About $90\%$ of errors lie inside; $5\%$ lie in each tail.
			
			\item \textbf{Packaging inspection ($\nu=5$).}
			\par Each excluded tail has probability $0.025$. Row $5$: $t^*=2.571$.
			\[L=-0.5-1.5(2.571)\approx\boxed{-4.36\text{ g}},\]
			\[U=-0.5+1.5(2.571)\approx\boxed{3.36\text{ g}}.\]
			About $95\%$ of mass errors lie between the limits, with $2.5\%$ on each side.
			
			\item \textbf{Processing-time alerts ($\nu=12$).}
			\par Lower tail $0.05$ gives $t_L=-1.782$; upper tail $0.10$ gives $t_U=1.356$ (row $12$).
			\[L=1+5(-1.782)\approx\boxed{-7.91\text{ minutes}},\]
			\[U=1+5(1.356)=\boxed{7.78\text{ minutes}}.\]
			About $5\%$ of outcomes lie below $L$, $10\%$ above $U$, and $85\%$ between them.
		\end{enumerate}
	\end{multicols}
\end{document}
